\documentclass[10pt,a4paper]{amsart}
\usepackage[margin=19mm]{geometry}
\usepackage{amsmath,amssymb,mathtools}
\usepackage[scr=rsfs,cal=euler]{mathalfa}
\usepackage{xfrac}
\usepackage[unicode,hidelinks,bookmarks=false]{hyperref}
\newcommand{\ie}{i.e.\ }
\newcommand{\cf}{cf.\ }
\newcommand{\resp}{respectively\ }
\newcommand{\Subsection}{\subsection}
\providecommand{\nobreakdash}{\nobreak\hbox{-}}
\setlength{\emergencystretch}{2em}
\multlinegap0pt
\usepackage[shortlabels]{enumitem}
\usepackage{amssymb}
\usepackage{float}
\newtheorem{definition}{Definition}
\newtheorem{theorem}{Theorem}
\newcommand{\calH}{\mathcal{H}}
\newcommand{\calV}{\mathcal{V}}
\title{Casorati Inequalities along Vertical and Horizontal Distributions of Pointwise Slant Riemannian Submersions}
\author{Tanveer Fatima, Sana Abdulkarem Alharbi,, Sharief Deshmukh}
\date{Source: arXiv:2608.15514v1 (CC BY 4.0)}
\begin{document}
\maketitle

\noindent\emph{Source and licence.} Casorati Inequalities along Vertical and Horizontal Distributions of Pointwise Slant Riemannian Submersions. Version arXiv:2608.15514v1 (CC BY 4.0), distributed by arXiv under the Creative Commons Attribution 4.0 International licence, http://creativecommons.org/licenses/by/4.0/, as recorded in the arXiv licence metadata for this exact version and shown on the abstract page https://arxiv.org/abs/2608.15514v1. Full licence notice: \texttt{licenses/arxiv-2608.15514-CC-BY-4.0.txt}.\par\smallskip

\noindent\textit{Editorial scope.} Counted unit: the article's theorem thm:sasakian\_subm\_horiz (source0.tex lines 1089--1100). The article's complete original proof of it is delivered with it (source0.tex lines 1101--1126, one \texttt{proof} environment of 26 lines). No mathematics is written, repaired or re-derived: the statement and the proof are the article's own lines, in the article's own notation, and no retained line is substituted or re-flowed.  Retained uncounted as the source-local context the proof needs: lines 358--364; lines 602--614; lines 729--742; lines 761--779; lines 841--852; lines 1003--1011.  Nothing was omitted from any retained span, so the package carries no omission record.  No retained line contains a citation, so the wrapper carries no bibliography and no bibliography entry.  No retained line contains a figure environment, an image-inclusion command or drawing code, so no image is delivered and the package declares no asset dependency.  Editorial formatting only, in addition: an AMS \texttt{amsart} preamble that restates the article's own declarations and macros used by the retained lines, namely the environments \texttt{definition}, \texttt{theorem} on separate counters, together with the standard packages those lines need.  The retained environments print in this document as Definition 1, Theorem 1 in order of appearance, whereas the article prints the same items with the numbering of its own full document.  The complete native source of the article as distributed in the arXiv e-print for this exact version is bundled unchanged as \texttt{sources/207754/source0.tex}.
\begin{align}\label{eq:gen_sasakian_curvature}
R(X,Y)Z &= c_1\{g(Y,Z)X - g(X,Z)Y\} \nonumber\\
&\quad + c_2\{g(X,\phi Z)\phi Y - g(Y,\phi Z)\phi X 
            + 2g(X,\phi Y)\phi Z\} \nonumber\\
&\quad + c_3\{\eta(X)\eta(Z)Y - \eta(Y)\eta(Z)X 
            + g(X,Z)\eta(Y)\xi - g(Y,Z)\eta(X)\xi\}
\end{align}

\begin{equation}\label{eq:norm_identity_contact}
  \|P^{\mathcal{V}}\|^2 =
  \begin{cases}
    (r-1)\cos^2\theta, & \xi_1\in\Gamma(\ker F_*),\\[2pt]
    r\cos^2\theta, & \xi_1\in\Gamma((\ker F_*)^{\perp}),
  \end{cases}
  \qquad
  \|P^{\mathcal{H}}\|^2 =
  \begin{cases}
    (s-1)\cos^2\theta, & \xi_1\in\Gamma((\ker F_*)^{\perp}),\\[2pt]
    s\cos^2\theta, & \xi_1\in\Gamma(\ker F_*).
  \end{cases}
\end{equation}

\begin{definition}\label{def:quasi-umbilical}
The fibers of a Riemannian submersion $F:(M_1,g_1)\to(M_2,g_2)$ are called
{invariantly quasi-umbilical} at $p\in M_1$ if there exist orthonormal
bases $\{U_1,\ldots,U_r\}$ of $\calV=\ker F_{*p}$ and
$\{U_{r+1},\ldots,U_{m_1}\}$ of $\calH=(\ker F_{*p})^{\perp}$ such that the
components $T_{ij}^{\alpha}=g_1(T_{U_i}U_j,U_\alpha)$ of O'Neill's tensor
satisfy
\begin{enumerate}[label=(\roman*)]
\item $T_{ij}^\alpha = 0$ for all $i \neq j$ and all
      $\alpha = r+1, \ldots, m_1$;
\item $T_{11}^\alpha = T_{22}^\alpha = \cdots =
      T_{r-1,r-1}^\alpha$ for all $\alpha = r+1, \ldots, m_1$.
\end{enumerate}
\end{definition}

\begin{theorem}\label{thm:complex_subm_vert}
Let $F:(M_1^{m_1}(c_1,c_2),g_1,J_1)\to (M_2^{m_2},g_2)$ be a pointwise slant Riemannian submersion from a generalized complex space form onto a Riemannian manifold with vertical space of dimension $r\ge 3$ and slant function $\theta$. Then
\begin{equation}\label{eq:complex_subm_vert_delta}
\rho^{\calV} \le \delta_C^{\calV}(r-1) + c_1 + \dfrac{3c_2}{r(r-1)}\|P^{\mathcal V}\|^2 = \delta_C^{\calV}(r-1) + c_1 + \dfrac{3c_2}{r-1}\cos^2\theta,
\end{equation}
and
\begin{equation}\label{eq:complex_subm_vert_delta_hat}
\rho^{\calV} \le \hat{\delta}_C^{\calV}(r-1) + c_1 + \dfrac{3c_2}{r(r-1)}\|P^{\mathcal V}\|^2 = \hat{\delta}_C^{\calV}(r-1) + c_1 + \dfrac{3c_2}{r-1}\cos^2\theta.
\end{equation}
Equality holds at $p \in M_1$ if and only if for suitable orthonormal 
bases, the components of the $T$-tensor satisfy
\[
T_{11}^{\alpha} = T_{22}^{\alpha} = \cdots = T_{r-1,r-1}^{\alpha} 
= \dfrac{1}{2}T_{rr}^{\alpha}, \quad T_{ij}^{\alpha} = 0 \text{ for } i\neq j,
\]
for all $\alpha = r+1,\ldots,m_1$. Geometrically, this means that the 
fibers of $F$ are invariantly quasi-umbilical in the sense of 
Definition~\ref{def:quasi-umbilical}. 
\end{theorem}

\begin{theorem}\label{thm:sasakian_subm_vert}
Let $F:(M_1^{m_1}(c_1,c_2,c_3),\phi_1,\xi_1,\eta_1,g_1)\to (M_2^{m_2},g_2)$ be a pointwise slant Riemannian submersion from a generalized Sasakian space form onto a Riemannian manifold with vertical space of dimension $r\ge 3$ and slant function $\theta$. Then
\begin{equation}\label{eq:sasakian_subm_vert}
\rho^{\calV} \le \begin{cases}
\delta_C^{\calV}(r-1) + c_1 + \dfrac{3c_2}{r}\cos^2\theta - \dfrac{2}{r}c_3, & \text{if } \xi_1 \in \Gamma(\ker F_*),\\[6pt]
\delta_C^{\calV}(r-1) + c_1 + \dfrac{3c_2}{r-1}\cos^2\theta, & \text{if } \xi_1 \in \Gamma(\ker F_*)^\perp,
\end{cases}
\end{equation}
and similarly for $\hat{\delta}_C^{\calV}$. Equality holds if and only if the fibers are invariantly quasi-umbilical 
 which is the same as in 
Theorem~\ref{thm:complex_subm_vert}.
\end{theorem}

\begin{theorem}\label{thm:complex_subm_horiz}
Let $F:(M_1^{m_1}(c_1,c_2),g_1,J_1)\to (M_2^{m_2},g_2)$ be a pointwise slant Riemannian submersion from a generalized complex space form onto a Riemannian manifold with horizontal space of dimension $s\ge 3$ and slant function $\theta$. Then
\begin{equation}\label{eq:complex_horiz}
\rho_{\calH}^{\calH} \le \delta_C^{\calH}(s-1) + c_1 + \dfrac{3c_2}{s-1}\cos^2\theta,
\end{equation}
and similarly for $\hat{\delta}_C^{\calH}$.Equality holds if and only if the tensor $A$ vanishes identically, i.e., 
$A_XY = 0$ for all $X,Y \in \Gamma((\ker F_*)^\perp)$, which is equivalent 
to the horizontal distribution being completely integrable.
\end{theorem}

\begin{theorem}\label{thm:sasakian_subm_horiz}
Let $F:(M_1^{m_1}(c_1,c_2,c_3),\phi_1,\xi_1,\eta_1,g_1)\to (M_2^{m_2},g_2)$ be a pointwise slant Riemannian submersion from a generalized Sasakian space form onto a Riemannian manifold with horizontal space of dimension $s\ge 3$ and slant function $\theta$. Then
\begin{equation}\label{eq:sasakian_horiz}
\rho_{\calH}^{\calH} \le \begin{cases}
\delta_C^{\calH}(s-1) + c_1 + \dfrac{3c_2}{s}\cos^2\theta - \dfrac{2}{s}c_3, & \text{if } \xi_1 \in \Gamma((\ker F_*)^\perp),\\[6pt]
\delta_C^{\calH}(s-1) + c_1 + \dfrac{3c_2}{s-1}\cos^2\theta, & \text{if } \xi_1 \in \Gamma(\ker F_*),
\end{cases}
\end{equation}
and similarly for $\hat{\delta}_C^{\calH}$. Equality holds if and only if the tensor $A$ vanishes identically, i.e., 
$A_XY = 0$ for all $X,Y \in \Gamma((\ker F_*)^\perp)$, which is equivalent 
to the horizontal distribution being completely integrable.
\end{theorem}

\begin{proof}
The argument parallels Theorem~\ref{thm:complex_subm_horiz} for the
$c_1,c_2$ terms and the optimization of the $A$-tensor, with the
additional $c_3$ (Reeb) term in~\eqref{eq:gen_sasakian_curvature}
contributing exactly as in Theorem~\ref{thm:sasakian_subm_vert},
according to whether $\xi_1$ is horizontal or vertical. By
\eqref{eq:norm_identity_contact},
$\|P^{\calH}\|^2=(s-1)\cos^2\theta$ when $\xi_1\in\Gamma((\ker F_*)^{\perp})$
and $\|P^{\calH}\|^2=s\cos^2\theta$ when $\xi_1\in\Gamma(\ker F_*)$, so
\[
\rho^{\calH}
= c_1+\frac{3c_2}{s(s-1)}\,(s-1)\cos^2\theta-\frac{2c_3}{s}
= c_1+\frac{3c_2}{s}\cos^2\theta-\frac{2c_3}{s},
\qquad \xi_1\in\Gamma((\ker F_*)^{\perp}),
\]
\[
\rho^{\calH}
= c_1+\frac{3c_2}{s(s-1)}\,s\cos^2\theta
= c_1+\frac{3c_2}{s-1}\cos^2\theta,
\qquad \xi_1\in\Gamma(\ker F_*).
\]
Substituting into the horizontal Casorati estimate
$\rho_{\calH}^{\calH}\leq\delta_C^{\calH}(s-1)+\rho^{\calH}$ obtained in
Theorem~\ref{thm:complex_subm_horiz} yields the two cases; the
$A\equiv0$ equality condition is unchanged from that theorem.
\end{proof}

\end{document}
